\section{Management Process}

\subsection{Project Estimates}

% Intentionally left blank for PA1 as instructed by the assignment brief.

\subsection{Project Plan}

The project follows an iterative, risk-driven process. Each sprint begins with
planning and ends with a review and retrospective. Safety-critical assumptions
and the highest technical risks are addressed before secondary features. The
dates below form an initial working schedule anchored to the PA1 deadline of
18 October 2026. They must be aligned with the official course calendar if the
teaching staff announces different milestones.

\subsubsection{Phase and Iteration Plan}

\begin{longtable}{|L{0.14\textwidth}|L{0.10\textwidth}|L{0.17\textwidth}|L{0.23\textwidth}|L{0.24\textwidth}|}
\hline
\textbf{Phase} & \textbf{Sprint} & \textbf{Dates} & \textbf{Objective} & \textbf{Milestone / exit criteria}\\ \hline
Inception & Sprint 1 & 5--18 Oct 2026 & Agree on the problem, users, scope, team process, major risks, initial plan, and product vision. & PA1 package completed; tracking tools active; 3--6 main UI screens and two weekly reports available.\\ \hline
Elaboration & Sprint 2 & 19 Oct--1 Nov 2026 & Baseline the critical requirements, architecture, data strategy, safety states, and principal user journey. & Reviewed backlog and use cases; architecture/UX baseline; data-source feasibility recorded.\\ \hline
Elaboration & Sprint 3 & 2--15 Nov 2026 & Retire the highest OCR, speech-to-text, NLP, normalization, interaction-data, and safety-communication risks through an executable proof of concept. & OCR, voice, and typed input converge on verified text; medication identity and a controlled interaction check are demonstrated.\\ \hline
Construction & Sprint 4 & 16--29 Nov 2026 & Implement and integrate the primary end-to-end medication safety workflow. & Test build supports intake through guidance; critical path passes integration tests.\\ \hline
Construction & Sprint 5 & 30 Nov--13 Dec 2026 & Complete prioritized reminder/follow-up capabilities, stabilize the prototype, and prepare the final demonstration. & Final prototype and documentation accepted; critical defects closed; demo and evidence ready.\\ \hline
\end{longtable}

\paragraph{Iteration objectives.}

\begin{longtable}{|L{0.11\textwidth}|L{0.28\textwidth}|L{0.38\textwidth}|L{0.13\textwidth}|}
\hline
\textbf{Sprint} & \textbf{Objectives} & \textbf{Key work items} & \textbf{Owner(s)}\\ \hline
Sprint 1 & Validate the project direction and establish a controllable project. & Refine proposal; draft Project Plan and Vision; create UI sketches; set up repository and tracker; define backlog, risks, working agreements, and weekly reports. & Long, Trang, Thành\\ \hline
Sprint 2 & Define what will be built and how the safety-critical workflow will work. & Analyze interview needs; write use cases and NFRs; prototype OCR, manual, typed, and voice flows; define NLP, OCR, speech-to-text, data, privacy, and safety interfaces and acceptance tests. & Long, Trang, Thành\\ \hline
Sprint 3 & Prove the highest-risk technical and safety assumptions. & Long implements NLP intent/entity processing and prepares, cleans, versions, and maps DAV/DDInter data; Trang implements OCR and image-based prescription intake; Thành implements voice capture and speech-to-text. Integrate all modes through editable verification and explicit assessment states. & Long, Trang, Thành\\ \hline
Sprint 4 & Deliver an integrated core product. & Implement primary pages and APIs; persist verified medication data; run automatic checks; show severity and limitations; simulate professional review; conduct system and usability testing. & Whole team\\ \hline
Sprint 5 & Complete, stabilize, and demonstrate the prioritized release. & Add confirmed-schedule reminders, severe-interaction gate, adherence and possible-ADR records; fix defects; run regression and acceptance tests; finalize documents and demo. & Whole team\\ \hline
\end{longtable}

\subsubsection{Builds and Releases}

Four executable builds are planned. Additional internal builds may be produced
by continuous integration, but the following are the review and test baselines.

\begin{longtable}{|L{0.16\textwidth}|L{0.16\textwidth}|L{0.18\textwidth}|L{0.40\textwidth}|}
\hline
\textbf{Build / release} & \textbf{Type} & \textbf{Target date} & \textbf{Purpose and included scope}\\ \hline
UI Prototype & Design prototype & 18 Oct 2026 & Validate navigation and OCR, manual, typed, and voice input; include transcript/extraction confirmation, results, warnings, and reminders; no clinical conclusions are produced.\\ \hline
Build 0.1 & Technical proof of concept & 15 Nov 2026 & Test OCR, voice-to-text, NLP, manual input, editable confirmation, DAV-based product identification, DDInter mapping, controlled interaction lookup, and explicit assessment states.\\ \hline
Build 0.2 & Integrated test build & 29 Nov 2026 & Provide the end-to-end core workflow for team, usability, integration, and safety testing.\\ \hline
Build 1.0 & Final course prototype & 13 Dec 2026 & Add prioritized reminder and follow-up functions, resolve critical defects, and provide a stable demonstration baseline.\\ \hline
\end{longtable}

Each baseline is tagged in version control and accompanied by a short release
note containing included scope, known limitations, test status, and unresolved
risks. A build is not described as safe for clinical use.

\subsubsection{Project Schedule}

The detailed schedule is maintained in the team's task-tracking tool. The
following table is the initial high-level baseline; individual backlog items
will be assigned and updated during sprint planning.

% PA1 submission requirement: export the overall Jira schedule/roadmap and
% insert it in this subsection. The image must show the three phases, five
% sprints, key tasks, assignees, and dates. Example:
% \begin{figure}[H]
%   \centering
%   \includegraphics[width=\textwidth]{images/jira-project-schedule.png}
%   \caption{Overall project schedule exported from Jira}
%   \label{fig:jira-project-schedule}
% \end{figure}

\begin{longtable}{|L{0.13\textwidth}|L{0.19\textwidth}|L{0.29\textwidth}|L{0.16\textwidth}|L{0.13\textwidth}|}
\hline
\textbf{Period} & \textbf{Workstream} & \textbf{Key scheduled work} & \textbf{Accountable owner} & \textbf{Output}\\ \hline
5--11 Oct & Product discovery & Consolidate proposal, problem statement, users, features, constraints, and initial risks. & Long & Scope baseline\\ \hline
12--18 Oct & PA1 planning & Complete plan/vision, UI prototype, tracker setup, weekly reports, and submission review. & Whole team & PA1 package\\ \hline
19--25 Oct & Requirements and UX & Interview synthesis; OCR, voice, typed, and manual-input use cases; editable verification; acceptance criteria; terminology. & Long / Trang / Thành & Requirements draft\\ \hline
26 Oct--1 Nov & Architecture and data preparation & Long prepares the DAV dataset, active-ingredient normalization rules, DDInter mapping approach, provenance, and coverage notes; the whole team defines multimodal interfaces, privacy boundaries, and the data model. & Long (data); whole team (architecture) & Baseline package\\ \hline
2--8 Nov & Multimodal intake spike & NLP and intent processing (Long); OCR and prescription extraction (Trang); voice capture and speech-to-text (Thành); shared verification contract. & Long / Trang / Thành & Working input slices\\ \hline
9--15 Nov & Intake and assessment integration & Long integrates the prepared DAV dataset and reviewed DDInter mappings; the team connects verified text from all input modes and tests severity, unresolved mappings, incomplete states, and explanation constraints. & Long leads; Trang and Thành support & Build 0.1\\ \hline
16--22 Nov & Core implementation & Implement and integrate intake, medication list, automatic checks, results, warnings, and limitations. & Whole team & Integrated alpha\\ \hline
23--29 Nov & Core verification & Professional-review simulation, system tests, usability checks, defect correction, release notes. & Whole team & Build 0.2\\ \hline
30 Nov--6 Dec & Follow-up features & Confirmed-schedule reminders, safety gate, adherence and possible-ADR records. & Trang / Thành & Feature-complete beta\\ \hline
7--13 Dec & Stabilization and handoff & Regression, acceptance review, critical defect closure, documentation, demo rehearsal, final baseline. & Whole team & Build 1.0\\ \hline
\end{longtable}

Dependencies are managed in this order: OCR extraction, speech transcription,
manual entry, and typed conversation each produce editable text; NLP processes
the applicable text; user verification precedes identification against DAV
data; Vietnamese product identification precedes active-ingredient mapping to
DDInter 2.0; successful mapping precedes interaction assessment; and assessment
status and severity precede guidance and reminder activation.
Documentation, risk review, testing, and tracking continue throughout all
sprints.

\subsubsection{Project Resourcing}

\begin{longtable}{|L{0.17\textwidth}|L{0.22\textwidth}|L{0.27\textwidth}|L{0.24\textwidth}|}
\hline
\textbf{Phase / sprint} & \textbf{People / roles} & \textbf{Skills and tools} & \textbf{Training or resource need}\\ \hline
Inception / Sprint 1 & Three team members; team leader coordinates. & LaTeX/Docs, Figma or equivalent, Jira or equivalent tracker, Git repository. & Agree on workflow, definition of done, medical-safety language, and document access.\\ \hline
Elaboration / Sprint 2 & Long leads NLP and medication-data preparation; Trang leads OCR; Thành leads speech-to-text; all review requirements and safety. & NLP evaluation, OCR evaluation, Web audio/microphone flows, responsive UX, architecture, DAV data preparation, and DDInter mapping analysis. & Long studies source fields, provenance, coverage, and terminology; the team defines consent, audio privacy, the verified-text schema, and representative scenarios.\\ \hline
Elaboration / Sprint 3 & Long owns NLP and all prepared DAV/DDInter data artifacts; Trang owns OCR; Thành owns voice-to-text; all integrate and test. & Language/intent tests, data cleaning and versioning, active-ingredient normalization and mapping, image preprocessing/OCR, speech recognition, and automated integration tests. & Long creates controlled medication and mapping datasets; Trang and Thành provide representative image and audio inputs for integration tests.\\ \hline
Construction / Sprint 4 & All developers and testers. & Web implementation, integration, browser testing, defect tracking, peer review. & Cross-train members on the integrated flow and failure-state testing.\\ \hline
Construction / Sprint 5 & All members; leader coordinates release. & Reminder logic, regression/acceptance testing, release management, presentation. & Long prepares the fallback demonstration dataset; the whole team verifies submission links and artifacts.\\ \hline
\end{longtable}

The project will use student laptops, modern desktop and mobile browsers, a
shared task tracker, a Git-based repository, a shared document workspace, and
approved external APIs or controlled datasets. No real patient data is needed
for development or demonstration.

\subsection{Project Monitoring and Control}

\subsubsection{Requirements Management}

The Vision document is the source of high-level stakeholder needs and product
features. Detailed requirements are represented as backlog items and use cases
in the task tracker. Each requirement or story receives a stable identifier,
priority, accountable owner, acceptance criteria, and links to related design,
implementation, and test evidence. Safety-critical items are labeled and must
include negative and incomplete-data acceptance criteria.

A proposed requirement change is recorded as a change item rather than applied
informally. The team analyzes its user value, safety impact, dependencies,
effort, schedule effect, and effect on committed sprint work. The Project
Manager and the member responsible for the affected area review the proposal;
changes to scope or safety behavior require agreement of the whole team. An
accepted change updates the Vision document or applicable specification, the
backlog and priority, acceptance tests, and traceability links. Rejected or
deferred changes remain recorded with their rationale.

Requirements are reviewed during sprint planning, clarified during the sprint,
demonstrated at the sprint review, and reconsidered during the retrospective.
No requirement is considered complete until its acceptance criteria and linked
tests pass and its user-facing limitations are documented.

\subsubsection{Reporting and Measurement}

% Intentionally left blank for PA1 as explicitly required by the brief.

\subsubsection{Risk Management}

Risks are identified from interviews, requirements, architecture work,
external dependencies, implementation, and testing. The team assesses each
risk by likelihood and impact and assigns an owner. The risk register is
reviewed at least during weekly coordination and at each sprint planning and
review meeting. A risk is closed only when its trigger can no longer occur or
the mitigation has been verified; realized risks become tracked issues.

\begin{longtable}{|L{0.07\textwidth}|L{0.10\textwidth}|L{0.27\textwidth}|L{0.25\textwidth}|L{0.12\textwidth}|L{0.09\textwidth}|}
\hline
\textbf{ID} & \textbf{Ranking} & \textbf{Description and impact} & \textbf{Mitigation / contingency} & \textbf{Owner} & \textbf{Status}\\ \hline
R-01 & High & A Vietnamese product or active ingredient may be identified incorrectly from DAV data or mapped to the wrong DDInter 2.0 entity, producing a misleading interaction result. & Require user verification; retain DAV source identifiers and mapping provenance; use reviewed mapping rules; test aliases and ambiguous ingredients; keep unresolved mappings unassessed. & Long & Open\\ \hline
R-02 & High & OCR errors could change a medication name, strength, or schedule and lead to an invalid assessment. & Display extracted fields for confirmation; show confidence/validation errors; provide manual correction; never assess unverified extraction. & Trang & Open\\ \hline
R-03 & High & AI-generated explanations could add unsupported medical conclusions or false reassurance. & Generate explanations only from structured findings; constrain prompts/output; prohibit independent safety conclusions; add disclaimers and professional escalation; test unsafe prompts. & Long & Open\\ \hline
R-04 & High & Speech-to-text may confuse Vietnamese accents, noisy speech, or medication terminology, causing an incorrect medication or schedule. & Show an editable transcript; require confirmation; highlight low-confidence terms; offer manual input; test accents, noise, dosages, and similar medication names; never assess unverified speech. & Thành & Open\\ \hline
R-05 & High & DAV publication limitations, DDInter coverage/access, or incompatibility between the two data sources may prevent reliable mapping and assessment. & Verify both sources in Elaboration; document dates and coverage; store only permitted data; maintain a reviewed mapping table; use a controlled demonstrator dataset or narrow scope if needed. & Long & Open\\ \hline
R-06 & High & Users may interpret ``no known interaction found'' as a guarantee of safety. & Use distinct result states and cautious wording; always display data limits; usability-test comprehension; route uncertainty for review. & Trang / Long & Open\\ \hline
R-07 & Medium & Sensitive prescription images, voice recordings, transcripts, or health data may be exposed through logs, repositories, or external services. & Obtain microphone permission; minimize or avoid audio retention; use synthetic/de-identified test data; restrict access; exclude secrets/data from Git; review evidence before publication. & Whole team & Open\\ \hline
R-08 & Medium & Severe-interaction reminder rules may be bypassed or may incorrectly block unrelated medication. & Implement and test the safety gate per medication; require review/confirmation for affected new reminders; preserve explicit status and audit evidence. & Thành & Open\\ \hline
R-09 & Medium & Three-person capacity or uneven workload may delay the five-sprint plan. & Keep a prioritized MVP; assign one owner per item; limit work in progress; review blockers weekly; cross-review and redistribute work at mid-sprint. & Trang & Open\\ \hline
R-10 & Medium & Integration may occur too late, causing incompatible NLP, OCR, speech-to-text, UI, service, or data contracts. & Define a common verified-text contract in Sprint 2; deliver thin slices in Sprint 3; integrate continuously; maintain automated smoke tests and tagged baselines. & Whole team & Open\\ \hline
R-11 & Medium & The lack of real professional integration may make review behavior unrealistic. & Mark it as a simulation; validate flows with available healthcare feedback; avoid claims of clinical deployment readiness. & Long / Trang & Open\\ \hline
\end{longtable}

\subsubsection{Configuration Management}

% Intentionally left blank for PA1; it will be completed in a later revision.

